
% ToC styling
\setlength{\cftbeforechapskip}{2pt}
\renewcommand{\cftchappresnum}{BAB }            % Set agar chapter punya awalan "BAB"
\renewcommand{\cftchapnumwidth}{3em}            % Tambah indentation agar tidak tabrakan
\renewcommand{\cftsecindent}{3em}               % Tambah indentation agar tidak tabrakan
\renewcommand{\cftchapfont}{\normalfont}        % Set agar chapter di daftar isi tidak bold
\renewcommand{\cftchappagefont}{\normalfont}    % Set agar numbering chapter tidak bold
\renewcommand{\cftdotsep}{0.75}                 % Set agar titik-titik tidak terlalu renggang
\renewcommand{\cftchapdotsep}{\cftdotsep}       % Set agar chapter juga punya titik-titik

\renewcommand{\cftfigindent}{0em}               % Set lof indent
\renewcommand{\cftfigpresnum}{Gambar }          % Set lof leading
\renewcommand{\cftfignumwidth}{5.5em}           % Set leading width\


% Daftar Hadir Bimbingan Table
\newenvironment{dhbtable}
{
    \begin{longtable}{| >{\centering\arraybackslash}m{1cm} | >{\justifying\arraybackslash}m{7cm} | >{\centering\arraybackslash}m{3cm} | >{\centering\arraybackslash}m{2cm} |} \hline
        \rule{0pt}{3ex} \textbf{No} &
        \rule{0pt}{3ex}\centering\textbf{Materi Bimbingan} &
        \rule{0pt}{3ex}\textbf{Tgl. Bimb.} &
        \rule{0pt}{3ex}\textbf{Paraf} \\ \hline
        \endfirsthead

        \hline
        \rule{0pt}{3ex}\textbf{No} &
        \textbf{Materi Bimbingan} &
        \textbf{Tgl. Bimb.} &
        \textbf{Paraf} 
        \endhead
}
{
    \end{longtable}
    \begin{flushright}
    {\city}, \textcolor{red}{\signaturedate}\\
    Mengetahui,\\
    Pembimbing {\jenisMBKM},\\[2cm]
    {\dosen}\\
\end{flushright}
}

% For dhbtable (dhbentry) -> insert new row
\newcommand{\dhbentry}[4]{
    #1 & 
    \fontsize{11}{14}\selectfont{#2} & 
    #3 & 
    #4
\\ \hline
}